\par\Needspace{10\baselineskip}\medskip
\hypertarget{native-128-article095}{}\pdfbookmark[2]{95.}{article-095}
\begin{center}95.\end{center}
We turn to the second task, determining the elements once the angle $g$ has been found. The semimajor axis is obtained immediately from formulas 10 and $10^{\ast}$. The following formulas may also be used in their place.
\[
\begin{aligned}
\mathrm{17)}\qquad a&=\dfrac{2mm\cos f\mathpunct{.}\BareRoot{rr'}}{yy\sin g^2}=\dfrac{kktt}{4yyrr'\cos f^2\sin g^2}\\[6pt]
\mathrm{17}^{\ast}\mathrm{)}\qquad a&=\dfrac{-2MM\cos f\mathpunct{.}\BareRoot{rr'}}{YY\sin g^2}=\dfrac{kktt}{4YYrr'\cos f^2\sin g^2}\mathpunct{\cdot}
\end{aligned}
\]
The semiminor axis $b=\BareRoot{ap}$ is obtained from equation 1, which, when combined with the preceding equations, gives
\[
\begin{aligned}
\mathrm{18)}\qquad p&=\left(\dfrac{yrr'\sin 2f}{kt}\right)^2\\[6pt]
\mathrm{18}^{\ast}\mathrm{)}\qquad p&=\left(\dfrac{Yrr'\sin 2f}{kt}\right)^2\mathpunct{\cdot}
\end{aligned}
\]
The area of the elliptic sector bounded by the two radius vectors and the elliptic arc is $=\frac12 kt\BareRoot{p}$, and the area of the triangle bounded by the same radius vectors and the chord is $=\frac12 rr'\sin 2f$. Hence the ratio of the sector to the triangle is $y:1$ or $Y:1$. This observation is of the greatest importance and also illustrates equations 12 and $12^{\ast}$ very neatly. It is clear from this that the terms $m$, $(l+x)^{\frac12}$, $X(l+x)^{\frac32}$ in equation 12 and in equation $12^{\ast}$ the terms $M$, $(L-x)^{\frac12}$, $X(L-x)^{\frac32}$ are proportional, respectively, to the area of the sector (between the radius vectors and the elliptic arc), the area of the triangle (between the radius vectors and the chord), and the area of the segment (between the arc and the chord). Clearly the first area equals either the sum or the difference of the other two, according as $v'-v$ lies between 0 and $180^0$ or between $180^0$ and $360^0$. If $v'-v$ is greater than $360^0$, the area of the entire ellipse must be counted in the area of the \SourceLeaf{129}sector and in that of the segment once for each complete revolution in the motion.

Since $b$ is $=a\cos\varphi$, combining equations 1, 10, and $10^{\ast}$ further gives
\[
\begin{aligned}
\mathrm{19)}\qquad \cos\varphi&=\dfrac{\sin g\operatorname{tang}f}{2(l+\sin\frac12 g^2)}\\[6pt]
\mathrm{19}^{\ast}\mathrm{)}\qquad \cos\varphi&=\dfrac{-\sin g\operatorname{tang}f}{2(L-\sin\frac12 g^2)}\mathpunct{\raisebox{.45ex}{,}}
\end{aligned}
\]
Substituting for $l$ and $L$ their values from Art. 89 then gives
\[
\mathrm{20)}\qquad \cos\varphi=\dfrac{\sin f\sin g}{1-\cos f\cos g+2\operatorname{tang}2\omega^2}\mathpunct{\cdot}
\]
This formula is unsuitable for the accurate calculation of the eccentricity when it is small. The following, more suitable formula is easily derived from it.
\[
\mathrm{21)}\qquad \operatorname{tang}\tfrac12\varphi^2=\dfrac{\sin\frac12(f-g)^2+\operatorname{tang}2\omega^2}{\sin\frac12(f+g)^2+\operatorname{tang}2\omega^2}\mathpunct{\raisebox{.45ex}{,}}
\]
It can also be given the following form by multiplying the numerator and denominator by $\cos 2\omega^2$.
\[
\mathrm{22)}\qquad \operatorname{tang}\tfrac12\varphi^2=\dfrac{\sin\frac12(f-g)^2+\cos\frac12(f-g)^2\sin 2\omega^2}{\sin\frac12(f+g)^2+\cos\frac12(f+g)^2\sin 2\omega^2}\mathpunct{\cdot}
\]
With either formula (using, if desired, auxiliary angles whose tangents are $\dfrac{\operatorname{tang}2\omega}{\sin\frac12(f-g)}\mathpunct{\raisebox{.45ex}{,}}$ $\dfrac{\operatorname{tang}2\omega}{\sin\frac12(f+g)}$ for the former, or $\dfrac{\sin 2\omega}{\operatorname{tang}\frac12(f-g)}\mathpunct{\raisebox{.45ex}{,}}$ $\dfrac{\sin 2\omega}{\operatorname{tang}\frac12(f+g)}$ for the latter), the angle $\varphi$ can always be determined with full precision.

The following formula may be used to determine the angle $G$. It follows directly by combining equations 5 and 7 with the unnumbered equation that follows the latter.
\[
\mathrm{23)}\qquad \operatorname{tang}G=\dfrac{(r'-r)\sin g}{(r'+r)\cos g-2\cos f\mathpunct{.}\BareRoot{rr'}}\mathpunct{\raisebox{.45ex}{,}}
\]
Introducing $\omega$ readily gives
\[
\mathrm{24)}\qquad \operatorname{tang}G=\dfrac{\sin g\sin 2\omega}{\cos 2\omega^2\sin\frac12(f-g)\sin\frac12(f+g)+\sin 2\omega^2\cos g}\mathpunct{\cdot}
\]
The remaining ambiguity is easily resolved with equation 7, which shows that $G$ must be taken between 0 and $180^0$ or between $180^0$ and $360^0$ according as the numerator in these two formulas is positive or negative.

\SourceLeaf{130}Combining equation 3 with the equations below, which come directly from equation II of Art. 8,
\[
\begin{aligned}
\dfrac1r-\dfrac1{r'}&=\dfrac{2e}{p}\sin f\sin F\\[6pt]
\dfrac1r+\dfrac1{r'}&=\dfrac2p+\dfrac{2e}{p}\cos f\cos F,
\end{aligned}
\]
readily gives the following formula.
\[
\mathrm{25)}\qquad \operatorname{tang}F=\dfrac{(r'-r)\sin f}{2\cos g\mathpunct{.}\BareRoot{rr'}-(r'+r)\cos f}\mathpunct{\raisebox{.45ex}{,}}
\]
Introducing the angle $\grave{\omega}$ into it gives
\[
\mathrm{26)}\qquad \operatorname{tang}F=\dfrac{\sin f\sin 2\omega}{\cos 2\omega^2\sin\frac12(f-g)\sin\frac12(f+g)-\sin 2\omega^2\cos f}\mathpunct{\cdot}
\]
The ambiguity is resolved here in the same way as before. — Once the angles $F$ and $G$ have been found, we have $v=F-f$, $v'=F+f$, from which the position of the perihelion will be \raisebox{.45ex}{.}known. We also have $E=G-g$, $E'=G+g$. Finally, the mean motion during the time $t$ is $=\dfrac{kt}{a^{\frac32}}=2g-2e\cos G\sin g$. The agreement of these expressions serves to check the calculation. The epoch of the mean anomaly corresponding to the instant midway between the two given instants is $G-e\sin G\cos g$ and may be transferred to any other time at will. It is somewhat more convenient still to compute the mean anomalies for the two given instants by the formulas $E-e\sin E$ and $E'-e\sin E'$, and to check the calculation by comparing their difference with $\dfrac{kt}{a^{\frac32}}$.
